%! Author = charon
%! Date = 6/4/24

\section{Verwandte Arbeiten}\label{sec:related-work}
Die Evaluierung und Analyse der Performance von Fuzzern ist ein Forschungsgebiet in der Cybersicherheitsforschung.
In dieser Sektion werden relevante wissenschaftliche Arbeiten und Methodiken zur Performance-Analyse von Netzwerkfuzzern
behandelt, insbesondere im Hinblick auf die Fuzzer Pulsar, AFLNet und boofuzz.\newline\newline
%In der wissenschaftlichen Arbeit von Wressnegger~\cite{pulsar} wird Pulsar als Fuzzer beschrieben, der ein feedback-orientiertes
%Modell verwendet, um systematisch Schwachstellen in Protokollen zu identifizieren.
%Die Forscher verwenden Benchmarking-Methoden, die auf der Analyse der Fehlerraten und der Geschwindigkeit der
%Fehlerentdeckung basieren.
%Eine entscheidende Methode zur Performance-Bewertung von Pulsar ist die Messung der Anzahl und Schwere der entdeckten
%Schwachstellen im Vergleich zu anderen Fuzzern.
%\gls{afl}Net, eine Erweiterung des bekannten American Fuzzy Lop für Netzwerkschnittstellen, wird in der Literatur~\cite{afl,AFLNet,AFLplusplus-Woot20,iot-fuzzing} als ein
%leistungsfähiges Werkzeug für das Fuzzing von Netzwerkprotokollen beschrieben.
%Die Performance von \gls{afl}Net wird durch Metriken wie die Codeabdeckung und die Anzahl entdeckter Schwachstellen gemessen.
%Die Autoren (u.a.\ Thorsten Holz und Moritz Schloegel) verwenden Benchmarking-Techniken, um die Effizienz von verschiedenen
%State of the Art Fuzzern bei der Erkennung von u.a.\ Protokollfehlern zu vergleichen und die Performance über verschiedene
%Testfälle hinweg zu bewerten~\cite{fuzzing-evaluation}.
%Boofuzz ist eine Weiterentwicklung des Sulley-Fuzzers und bietet eine umfassende Plattform für das Netzwerkfuzzing.
%Die Performance-Bewertung von boofuzz erfolgt durch die Analyse der Anzahl und Art der entdeckten Fehler.
%Die wissenschaftliche Literatur bietet verschiedene Ansätze zur Performance-Analyse von Netzwerkfuzzern.
%In der Arbeit von Holz werden zwei Hauptmethoden zur Performance-Bewertung vorgestellt: die quantitative und die qualitative Analyse.
%\begin{itemize}
%    \item Quantitative Analyse: Diese Methode umfasst die Messung von Metriken wie der Anzahl entdeckter Schwachstellen.
%    Quantitative Methoden sind in der Arbeit von detailliert beschrieben und beinhalten die Berechnung von Kennzahlen wie der Rate der Fehlerentdeckung und der Effizienz der Testausführung.
%    Die Forscher bewerten, wie schnell und effektiv verschiedene Fuzzer Schwachstellen finden und welche Menge an Testfällen benötigt wird, um signifikante Fehler zu identifizieren~\cite{fuzzing-evaluation}.
%    \item Qualitative Analyse: Diese Methode konzentriert sich auf die Qualität der gefundenen Schwachstellen und deren praktische Relevanz.
%    Es wird untersucht, wie die Reproduzierbarkeit und Ausnutzbarkeit der Schwachstellen bewertet werden können.
%    Die Forscher analysieren, ob die gesammelten Daten über mehrere Publikationen reproduzierbar sind oder ob die Entdeckung
%    von Schwachstellen auf einer Art von Zufall basieren~\cite{fuzzing-evaluation}.
%\end{itemize}
%Eine weitere Arbeit von Görz und Mathis zeigt eine weitere Benchmarking-Methode für Fuzzer, in der die Performance der
%Mutationen und Generatoren von Fuzzern analysiert~\cite{mutuation-analysis} wird.\newline
%Zusammenfassend zeigen die wissenschaftlichen Veröffentlichungen, dass die Performance von Netzwerkfuzzern wie Pulsar,
%\gls{afl}Net und boofuzz durch eine Kombination von quantitativen und qualitativen Methoden bewertet wird.
%Die Analyse der Performance erfolgt durch die Untersuchung von Metriken wie Fehlerentdeckungsrate, sowie durch die Bewertung
%der Qualität und Relevanz der gefundenen Schwachstellen.
%Die Literatur bietet umfassende Benchmarking-Methoden und Fallstudien, die eine detaillierte Bewertung der Fuzzer-Leistung
%ermöglichen und zur Weiterentwicklung der Fuzzing-Techniken beitragen.
%
- Fuzzing methodology ->  iot, A survey, evaluating directed fuzzers, SoK, a survey of network prot fuzz
- Fuzzing trends -> Pulsar, A survey, Protocol reversing, AFLNet,
- Fuzz benching -> SoK, Effective, Improving fuzzing assessment, Evaluating directed fuzzers,
- Network prot testing mit ML/AI -> CNN, Pulsar, SoK, A survey of network prot fuzz, Fuzzing state of the art, Learning Guided network fuzzing for testing cyber-physical systems
- Ausreiser: AI Based Framework for Automatic Test Data Generation, Bit-oriented format extraction approach for automatic binary protocol reverse engineering, IoTFuzzer: Discovering Memory Corruptions in IoT Through App-based Fuzzing | Welcome,
            Systematic Assessment of Fuzzers using Mutation Analysis, Towards Effective Performance Fuzzing

Beim Fuzzing von Applikationen entwickeln sich mit den sich ständig wechselnden technologischen Fortschritten auch die
Anforderungen und Herangehensweisen an das Fuzzing.
Mit einer immer weiter zunehmenden Anzahl an \gls{iot}-Geräten und der zunehmenden Vernetzung von Geräten und Systemen
werden auch die Anforderungen an das Fuzzing immer komplexer.
Somit ist es notwendig, dass Fuzzer nicht nur Applikationen, sondern auch Netzwerkprotokolle testen können.
Holz et al.~\cite{fuzzing-evaluation,mutuation-analysis} und Flores et al.~\cite{fuzzing-assessment} beschreiben in ihrer Arbeit die Notwendigkeit, dass Fuzzer auch mit sich ständig
ändernden Anforderungen eine Evaluation und Performance-Analyse von Fuzzern eine wichtige Rolle bei der Auswahl des richtigen
Fuzzers spielt.
Hierzu wurden Herangehensweisen und Methoden entwickelt, um die Performance von Fuzzern zu bewerten und diese Ansätze als
neuen Standard zu empfehlen.\newline
Desweiteren werden neuere Ansätze und Methoden zur Analyse von Netzwerkprotokollen vorgestellt.
Diese Ansätze basieren auf Machine Learning und künstlicher Intelligenz und ermöglichen es, Netzwerkprotokolle zu analysieren
und Schwachstellen zu identifizieren.
Der Ansatz von Gascon~\cite{pulsar} basiert auf der Verwendung einer Markov-Kette, um das Verhalten von Netzwerkprotokollen
zu modellieren und anhand des erlernten Zustandsmodells eingaben zu generieren und unerwartete Zustände zu provozieren.
Die Arbeit von Garshasbi~\cite{cnn} beschreibt die Verwendung von Convolutional Neural Networks, um Netzwerkprotokolle zu analysieren.
Hierbei wird ein Modell trainiert um Netzwerkprotokolle und vorallem die Struktur der Eingaben des Protokolls
auf binärer Ebene zu erlernen.
Das Erlernen der Struktur einer Eingabe und der damit verbundenen Zustandsmaschine eines Programms wird ebenfalls mit
Fuzzern wie AFLNet~\cite{AFLNet} mit der bereits bewährten Methode der Codepfadabdeckung durchgeführt.\newline
Die Auswertung der Güte von Fuzzern ist ein weiterer Teilbereich der Fuzzing-Forschung.
Hierbei werden bereits existierende Standards, wie Googles FuzzBench Suite~\cite{fuzzbench} oder LAVA~\cite{lava}, verwendet,
um die Performance von Fuzzern zu bewerten.

%TODO: Mehr ausführen und vorallem mehr auf IoT eingehen!
% Mehr auf benchmarking eingehen, wie SoK. Nimm dir die Arbeit ab und gehe auf diesen "Heiligen Graal" mehr ein.
